\lyseis{6.4}{Γραφική παράσταση της αναλογίας}
\sol{1} $y=0,\,2,\,4,\,6,\,8,\,10$. Τα σημεία $(0,\,0)$, $(1,\,2)$, \dots, $(5,\,10)$ βρίσκονται σε ημιευθεία από το $O$.
\sol{2} Και οι τρεις είναι ημιευθείες από το $O$. Η $y=2x$ είναι η πιο απότομη (μεγαλύτερος συντελεστής) και η $y=0{,}5x$ η πιο ήπια.
\sol{3} α) Ναι ($\alpha=1{,}5$). β) Όχι ($y=x+1$, δεν περνά από το $O$). γ) Όχι ($\frac{1}{1}=1$, $\frac{3}{2}=1{,}5$). δ) Ναι ($\alpha=2{,}5$).
\sol{4} Τα πρώτα τέσσερα: ναι, $\alpha=1{,}5$. Στα $(2,\,3)$, $(4,\,6)$, $(5,\,8)$ όχι, γιατί $\frac{8}{5}=1{,}6\neq1{,}5$.
\sol{5} $\alpha=\frac{12}{5}=2{,}4$. Για $x=2$: $y=4{,}8$. Για $y=30$: $x=30:2{,}4=12{,}5$.
\sol{6} Πίνακας $y=2,\,3,\,4,\,5$. Τα σημεία είναι σε ευθεία που τέμνει τον $Oy$ στο $2$ και δεν περνά από το $O$, άρα δεν είναι αναλογία (για $x=0$ είναι $y=2\neq0$).
\sol{7} $y=8x$. Για $2{,}5$ κιλά: $20$ €. Με $36$ €: $36:8=4{,}5$ κιλά.
\sol{8} Μετά από $3$ ώρες $54$ και $36$ km, διαφορά $18$ km. Ο Βασίλης κάνει $30$ km σε $30:12=2{,}5$ ώρες.
\sol{9} $y=1{,}85x$. Για $12$ λίτρα $22{,}2$ €.
\sol{10} $y=0{,}4x$. Για $3{,}5$ μ.: $1{,}4$ kg.
\sol{11} Γιατί όταν $x=0$ είναι και $y=0$: αν δεν αγοράσουμε τίποτα, δεν πληρώνουμε τίποτα. Άρα το $(0,\,0)$ είναι πάντα σημείο της.
\sol{12} Όχι. Για ανάλογα ποσά, όταν $x=0$ είναι $y=0$, άρα η γραφική παράσταση πρέπει να περνά από το $O$.
\sol{13} $\alpha=\frac{4}{3}$. Για $x=9$: $y=12$. Για $y=2$: $x=1{,}5$.
\sol{14} Η θερμοκρασία δεν είναι ανάλογη της ώρας: δεν είναι $0$ όταν η ώρα είναι $0$ και ο λόγος δεν μένει σταθερός.
\sol{15} Ελεύθερη απάντηση. Ένα βήμα είναι περίπου $0{,}6$--$0{,}7$ μ. και τα σημεία βρίσκονται περίπου σε ημιευθεία από το $O$.
\sol{16} Ελεύθερη απάντηση.
